\section{Feature engineering}
\label{sec:feature-engineering}

Table~\ref{tab:features} lists every declared feature, generated directly
from \texttt{conf/features.yaml} -- the single source of truth consumed by
both the training job and the serving path
(Section~\ref{sec:methodology-serving}).

% Generated by scripts/export_feature_tables.py -- do not edit by hand.
\begin{table}[t]
  \centering
  \caption{Engineered features. Static attributes need no cut-off;
  historical aggregates are computed strictly before the row being
  predicted (point-in-time).}
  \label{tab:features}
  \begin{tabular}{llll}
    \toprule
    Feature & Family & Type & Point-in-time \\
    \midrule
    \texttt{user\_height\_cm} & user & static & -- \\
    \texttt{user\_weight\_kg} & user & static & -- \\
    \texttt{user\_bmi} & user & static & -- \\
    \texttt{user\_bust\_band} & user & static & -- \\
    \texttt{user\_bust\_cup\_ordinal} & user & static & -- \\
    \texttt{user\_body\_type} & user & static & -- \\
    \texttt{user\_age} & user & static & -- \\
    \texttt{item\_category} & item & static & -- \\
    \texttt{item\_nominal\_size} & item & static & -- \\
    \texttt{user\_prior\_review\_count} & interaction & historical aggregate & \checkmark \\
    \texttt{user\_prior\_mismatch\_rate} & interaction & historical aggregate & \checkmark \\
    \texttt{user\_prior\_small\_rate} & interaction & historical aggregate & \checkmark \\
    \texttt{item\_prior\_mismatch\_rate} & interaction & historical aggregate & \checkmark \\
    \texttt{context\_rented\_for} & context & static & -- \\
    \bottomrule
  \end{tabular}
\end{table}


\subsection{Static versus historical aggregate: the distinction that matters}

This dataset makes a distinction that is easy to blur and important to get
right.

\textbf{Static attributes} -- height, weight, BMI, bust band and cup, body
type, age, item category -- are known the moment a user fills in a profile
or an item enters the catalogue. They carry \emph{no} point-in-time
cut-off because there is no "future" version of a person's height to leak:
these values simply exist before any review is written, train or test
alike. \texttt{recommender.features.user\_features} and
\texttt{item\_features.item\_category} read them directly off the row being
featurized.

\textbf{Historical aggregates} -- a user's or an item's past mismatch rate,
and a user's prior review count -- are a fundamentally different kind of
feature: they summarize \emph{outcomes of other reviews}, and the review
being predicted is one such outcome. If a user's "prior mismatch rate" were
computed over \emph{all} of that user's reviews (including the one being
predicted, or ones that happened later), the feature would partially encode
the label it is meant to help predict. This is exactly the leak this project
is built to avoid.

\subsection{Point-in-time construction}

Each review carries two chronological positions, assigned once when the
dataset is indexed (\texttt{recommender.data.loader.build\_indexed\_dataset},
\texttt{recommender.features.pipeline.assign\_global\_order}):

\begin{itemize}
  \item \texttt{review\_order} -- the review's position within its own
        user's history (0, 1, 2, \ldots), used by the two user-level
        aggregates.
  \item \texttt{global\_order} -- the review's position within the entire
        dataset's chronology (by \texttt{review\_date}), used by the
        item-level aggregate, since an item can be reviewed by many
        different users across every partition.
\end{itemize}

For a review at position $k$ (per the relevant order column), the historical
aggregate is an \emph{expanding} statistic over strictly prior positions
only:
\begin{equation}
  \widehat{r}(k) = \frac{1}{k} \sum_{j < k} \mathbb{1}[\text{outcome}_j
  \text{ was a mismatch}], \qquad \widehat{r}(0) \text{ undefined (no prior
  history)}.
\end{equation}
The strict inequality is the entire contract. A review at a given position
can see every review strictly before it, never the review at its own
position, and never one after it. This is enforced by computing the
statistic over a frame sorted by the order column and using an explicit
less-than comparison, not a window function that could silently include the
current row (\texttt{interaction\_features.\_expanding\_prior\_mean}).

Rows with no qualifying prior history (a user's or item's very first
appearance) get \texttt{NaN} from the raw computation; the feature pipeline
then applies the fill policy declared in \texttt{conf/features.yaml} -- the
global mismatch rate for the two rate features, zero for the count and the
small-rate feature, since a user's first review has, by definition, no
observed directional bias yet.

\subsection{Leakage tests}

\texttt{tests/test\_features\_leakage.py} operationalizes the point-in-time
contract directly: it injects a fabricated future review (an extreme, label-
flipping outcome) for an existing user or item, strictly after the real
cut-off, and asserts that every feature value for rows at or before that
cut-off is unchanged. A second test asserts the point-in-time aggregate
differs from a naive whole-history average whenever a user has more than one
review -- if they ever coincided, every row would be seeing its own future.

\subsection{Context features}

A shopper-stated occasion (\texttt{context\_rented\_for}, e.g.\ "wedding",
"work") is supplied at request time, like a search filter. It is information
about the current transaction, not a historical aggregate, so it needs no
cut-off either -- it exists in exactly the same sense static attributes do.

\subsection{Training/serving parity}

\texttt{FeatureStore.transform} (batch) and \texttt{FeatureStore.row\_vector}
(single serving-time row) are the only two entry points into the feature
pipeline, and \texttt{tests/test\_features\_parity.py} asserts they produce
numerically identical output for the same logical row -- the failure mode
this guards against is an offline batch computing one feature vector while
an online request computes a different one for what should be the same
input.
